\section{The literal canonical branch}
\label{sec:literal-canonical-branch}

The robustness theorem deliberately did not assume that \(b\) came
from the TPC-32 canonical common-target projector.  We now restore
that provenance exactly by importing the established TPC-93
Theorem 3.3, Corollary 3.4, and Theorem 3.5 (Complete Decorated
Affine Export), whose explicit inverse and multiplicity reconstruction
furnish the required lossless source-to-decorated ledger.

For a left-polarized TPC-93 key,
\[
 \sigma_\theta U_\theta(t)
 =
 M_\theta(t)j_\theta+h_0,
\qquad
 P_\theta=n_\theta j_\theta+h_0.
\]
An active content-inversion key satisfies
\[
 b=c\kappa\mid\sigma_\theta U_\theta(t),
 \qquad
 b\mid P_\theta.
\]
Thus \(b\) is literally a common divisor of the two original targets,
not merely a divisor of a reduced surrogate.
For a right-polarized key the two row labels are exchanged.  The
same common-target statements hold, while the determinant-character
orientation is recorded separately by
\(\varepsilon_\theta=-1\).

\begin{theorem}[Literal removal of the robustness content]
\label{thm:canonical-H-one}
On the TPC-32 primitive target support, import the established
TPC-93 Theorem 3.3, Corollary 3.4, and Theorem 3.5, including their
lossless inverse provenance and multiplicity reconstruction
\citep{WangTPC93}.  Then
every active content-inversion key in either polarization satisfies
\begin{equation}
 (b,h_0)=1.
\label{eq:canonical-b-coprime-h}
\end{equation}
Consequently
\[
 (B,h_0)=1,\qquad H=1.
\]
The two affine M\"obius arguments after extracting \(B\) are
individually primitive and have the same prescribed determinant
\(h_0\):
\[
 D_\theta(\tau+Bz),\qquad
 V_{\theta,b}(z)=\frac{U_\theta(\tau+Bz)}B.
\]
\end{theorem}

\begin{proof}
The TPC-32 common-target-divisor lemma states that every common
target divisor \(w\) is coprime to the physical product containing
the fixed shift \(h_0\) \citep{WangTPC32}.  Apply it with
\[
 w=b,\quad
 N_m=\sigma_\theta U_\theta(t),\quad
 N_n=P_\theta.
\]
This proves \eqref{eq:canonical-b-coprime-h}.  Since \(B\mid b\),
\((B,h_0)=1\).  Now apply
\cref{thm:robustness-content} to obtain \(H=1\) and the primitive
determinant \(h_0\).
\end{proof}

\begin{corollary}[Literal normalized M\"obius pair]
\label{cor:literal-normalized-mobius-pair}
On the literal canonical branch,
\begin{align}
 &\mu(D_\theta(\tau+Bz))
 \mu(U_\theta(\tau+Bz))
\notag\\
 &\quad=
 \mu(B)
 \mu(D_\theta(\tau+Bz))
 \mu(V_{\theta,b}(z))
 \one_{(B,V_{\theta,b}(z))=1}.
\label{eq:literal-normalized-mobius-pair}
\end{align}
Every squareful \(B\) branch is identically zero.  Every surviving
branch is a primitive determinant-\(h_0\) TPC-90 pair with one extra
coprimality mask and the extracted sign \(\mu(B)\)
\citep{WangTPC90}.
\end{corollary}

\begin{proof}
Set \(H=1\) in
\eqref{eq:two-stage-mobius-normalization}.
\end{proof}

\begin{proposition}[Literal content divisibility]
\label{prop:literal-content-divides-phase-slope}
Under the imported lossless-provenance theorem above, for every
active canonical content key in either polarization and every
\(z\in\Z\),
\begin{equation}
 b\mid M_\theta(\tau+Bz)-n_\theta.
\label{eq:literal-b-divides-row-difference}
\end{equation}
Consequently
\begin{equation}
 b\mid\Omega_{\theta,b}
 =\ell_\theta v_\theta\sigma_\theta B,
 \qquad
 c\mid\Omega_{\theta,b}.
\label{eq:literal-content-divides-phase-slope}
\end{equation}
\end{proposition}

\begin{proof}
The defining progression gives
\[
 b\mid M_\theta(\tau+Bz)j_\theta+h_0,
 \qquad
 b\mid n_\theta j_\theta+h_0
\]
for every integer \(z\).  Subtraction gives
\[
 b\mid
 \{M_\theta(\tau+Bz)-n_\theta\}j_\theta.
\]
The TPC-32 common-target-divisor lemma gives
\((b,j_\theta)=1\), so \(j_\theta\) may be cancelled, proving
\eqref{eq:literal-b-divides-row-difference}.  Subtract that
divisibility at \(z+1\) and \(z\).  The difference is
\(\Omega_{\theta,b}\), proving
\eqref{eq:literal-content-divides-phase-slope}.  Finally \(c\mid b\).
The right polarization is identical after exchanging the two row
labels; the sign \(\varepsilon_\theta=-1\) affects the oriented phase
but not any displayed divisibility.
\end{proof}

\begin{remark}[Why the abstract palette is still needed]
The conclusion \(H=1\) depends on three exact provenance facts:
\(\sigma U\) and \(P\) are the original targets, \(b\) divides both,
and the TPC-32 primitive support has not been enlarged.  If any of
these facts is lost under a future completion or surrogate, one must
return to \cref{thm:robustness-content}; the correct primitive
determinant is then \(h_0/H\), not a silently relabeled \(h_0\).
\end{remark}
